\section*{第 \ref{ch:b2:plane-waves-polarization} 章——平面电磁波与偏振}

\begin{solution}{exo:b2:plane-waves-polarization:1}
$\lambda = \qty{3.0}{m}$，$k = \qty{2.1}{rad/m}$，$T = \qty{10}{ns}$；$E_0 = \sqrt{2I/\varepsilon_0c}
= \qty{27}{mV/m}$，$B_0 = \qty{9e-11}{T}$；$u = I/c = \qty{3.3e-15}{J/m^3}$；电动势
$E_0 \times \qty{1}{m} = \qty{27}{mV}$；$I/hf = 10^{-6}/6.6 \times 10^{-26} = \qty{1.5e19}{m^{-2}.s^{-1}}$。
\end{solution}

\begin{solution}{exo:b2:plane-waves-polarization:2}
\qty{300}{kHz}（\qty{1.2e-9}{eV}，无线电波）；\qty{3}{GHz}（\qty{1.2e-5}{eV}，微波）；
\qty{30}{THz}（\qty{0.12}{eV}，红外线）；\qty{545}{THz}（\qty{2.25}{eV}，可见光）；
\qty{3e16}{Hz}（\qty{124}{eV}，极紫外）；\qty{3e18}{Hz}（\qty{12}{keV}，X 射线）；
\qty{3e20}{Hz}（\qty{1.2}{MeV}，γ 射线）。\qty{1}{MeV}：\qty{1.2}{pm}；\qty{160}{GHz}：\qty{1.9}{mm}。
\end{solution}

\begin{solution}{exo:b2:plane-waves-polarization:3}
$0.75I_0$，$0.56I_0$，$0.42I_0$；\omterm{def:b2:plane-waves-polarization:states}{自然光}：$0.375$，$0.28$，$0.21$。去掉
中间那块时：$0.75 \times \cos^260^\circ = 0.19I_0$（\omterm{def:b2:plane-waves-polarization:states}{自然光}为 $0.125I_0$）：
中间的起偏器通过
转动偏振方向，使\emph{更多}的光透过。
\end{solution}

\begin{solution}{exo:b2:plane-waves-polarization:4}
$\operatorname{div}\vect E = \partial_xE_x = 0$，$\operatorname{div}\vect B = 0$；
$\operatorname{\vect{curl}}\vect E = E_0k\sin(\omega t - kz)\,\vect e_y = -\partial_t\vect B$ 当且仅当 $\omega =
kc$——法拉第同时确定了方向（$\vect e_y$）与
大小（$E_0/c$）的 $\vect B$；$\operatorname{\vect{curl}}\vect B = -(E_0k/c)\sin(\omega t - kz)\,\vect e_x
= \mu_0\varepsilon_0\partial_t\vect E$。若 $\vect E$ 沿 $\vect e_z$，则 $\operatorname{div}
\vect E \ne 0$，而无电荷：不可能——这种波是横波。
\end{solution}

\begin{solution}{exo:b2:plane-waves-polarization:5}
(a) $\vect E = E_0(\cos(\omega t - kz), \pm\sin(\omega t - kz), 0)$：模为 $E_0$，
角度为 $\pm(\omega t - kz)$。(b) $(\cos, \sin) + (\cos, -\sin) = (2\cos, 0)$：\omterm{def:b2:plane-waves-polarization:states}{线
偏振}，沿 $x$ 方向；相位差 $\varphi$ 使方向转过 $\varphi/2$。(c)
在每个角度下均为 $I_0/2$。(d) \omterm{prop:b2:plane-waves-polarization:plates}{四分之一波片}将\omterm{def:b2:plane-waves-polarization:states}{圆偏振}光
变为\omterm{def:b2:plane-waves-polarization:states}{线偏振}光，随后起偏器将其消光；\omterm{def:b2:plane-waves-polarization:states}{自然光}
仍为非偏振光。
\end{solution}

\begin{solution}{exo:b2:plane-waves-polarization:6}
(a) $e = \lambda/4(n_s - n_f) = \qty{16}{\micro m}$；\omterm{prop:b2:plane-waves-polarization:plates}{半波片}厚度为 \qty{32}{\micro m}。(b) $\delta =
(\pi/2)(589/450) = \qty{2.1}{rad}$：\omterm{def:b2:plane-waves-polarization:states}{椭圆偏振}光。(c) $\delta \propto e/\lambda$：
厚度为 $m + \tfrac14$ 个波长的\omterm{prop:b2:plane-waves-polarization:plates}{波片}，其 $\dd\delta/\dd\lambda$ 与 $\dd\delta/\dd T$ 大了
$4m + 1$ 倍。(d) $\delta \propto e$：在 \qty{16}{\micro m} 上，一微米的误差对应的相位为 $6\%$ 的 $\pi/2$。
\end{solution}

\begin{solution}{exo:b2:plane-waves-polarization:7}
(a) 慢分量的符号改变：$(\cos\alpha, \sin\alpha) \to (\cos\alpha,
-\sin\alpha)$。(b) 偏振方向为 $\alpha$ 的光入射到方位为 $\beta$ 的\omterm{prop:b2:plane-waves-polarization:plates}{波片}上，
出射时变为 $2\beta - \alpha$：转角为 $2\beta$；即 \qty{45}{\degree}。(c)
反转手性。(d) $I = I_0\sin^22\beta$：每转
一周有四次极大与四次消光。
\end{solution}

\begin{solution}{exo:b2:plane-waves-polarization:8}
$\cos^{2N}(\pi/2N)$：$0$，$0.25$，$0.61$，$0.78$，$0.976$，$\to 1$。
连续扭转的介质（旋光现象、液晶盒）
无损耗地转动偏振方向。
\end{solution}

\begin{solution}{exo:b2:plane-waves-polarization:9}
(a) $I_{\max} = I_n/2 + I_p$，$I_{\min} = I_n/2$；$p = I_p/(I_n + I_p) = (I_{\max} -
I_{\min})/(I_{\max} + I_{\min})$。(b) $p = 4/6 = 0.67$。(c) 与之正交的起偏器
透过 $I_n/2 = 0.05I$：$95\%$ 被滤除。(d) 不能——
两者给出的都是恒定强度；需要\omterm{prop:b2:plane-waves-polarization:plates}{四分之一波片}。
\end{solution}

\begin{solution}{exo:b2:plane-waves-polarization:10}
(a) 相位 $\omega t - k(z\cos\theta \pm x\sin\theta)$：求和得 $2E_0\cos(kx\sin\theta)\cos(\omega t
- kz\cos\theta)\,\vect e_y$。(b) $\Delta x = \lambda/2\sin\theta$：在 \qty{5}{\degree} 时为 \qty{3.6}{\micro m}，
在 \qty{90}{\degree} 时为 $\lambda/2 = \qty{316}{nm}$。(c) $v_\varphi = c/\cos\theta > c$：这是图样的
\omterm{def:b2:dispersion-wave-packets:relation}{相速度}；能量沿 $z$ 方向以 $c\cos\theta$ 传播（每列波以 $c$ 传播）。
(d) $B_z = (\sin\theta/c)(E_1 - E_2)$，$E_y = E_1 + E_2$：$\langle\Pi_x\rangle \propto
\langle E_1^2 - E_2^2\rangle = 0$；$\Pi_z \propto \cos\theta(E_1 + E_2)^2$，在
亮面处最大。
\end{solution}

\begin{solution}{exo:b2:plane-waves-polarization:11}
(a) 无电压时：旋转 \qty{90}{\degree}，通过正交起偏器：
亮；有电压时：暗。(b) 漏光 $\sin^25^\circ = 7.6 \times 10^{-3}$：对比度
$130$。(c) 扭转结构相当于一叠延迟器，其 $\delta \propto 1/\lambda$：
在 \qty{550}{nm} 处精确，而光谱两端的旋转有偏差，
从而漏光——暗态看起来带紫色。(d) 当外部起偏器
与屏幕的输出起偏器正交时；该角度即
屏幕的偏振轴。
\end{solution}

\begin{solution}{exo:b2:plane-waves-polarization:12}
(a) $\vect e_x = \tfrac12[(\vect e_x + \iu\vect e_y) + (\vect e_x - \iu\vect e_y)]$；经过 $\ell$ 后
两个分量的相位分别为 $k_L\ell$ 与 $k_R\ell$，重新合成为
\omterm{def:b2:plane-waves-polarization:states}{线偏振}光，其偏振方向角为 $(k_L - k_R)\ell/2 = \pi(n_L - n_R)\ell/\lambda$。(b) $n_L -
n_R = \theta\lambda/\pi\ell = 0.116 \times 589 \times 10^{-9}/0.314 = \qty{2.2e-7}{}$。(c)
$0.01/6.65 \times 0.1 = \qty{1.5e-4}{g/mL}$。(d) 法拉第旋转由
$\vect B$ 决定，与传播方向无关：返程时它同向叠加
（\qty{45}{\degree} $+$ \qty{45}{\degree} $=$ \qty{90}{\degree}），于是输入起偏器挡住
返回的光——这就是隔离器。
\end{solution}

\begin{solution}{pb:b2:plane-waves-polarization:1}
\textbf{1.} $f = \qty{4.74e14}{Hz}$，$k = \qty{9.9e6}{rad/m}$，$hf = \qty{3.1e-19}{J} =
\qty{1.96}{eV}$；每秒 $3.2 \times 10^{16}$ 个光子。

\textbf{2.} $I = 0.01/\pi(0.5 \times 10^{-3})^2 = \qty{1.3e4}{W/m^2}$；$E_0 = \qty{3.1}{kV/m}$，
$B_0 = \qty{10}{\micro T}$；$\vect E = E_0\cos(\omega t - kz)\vect e_x$，$\vect B = (E_0/c)\cos(\omega t
- kz)\vect e_y$。

\textbf{3.} $u = I/c = \qty{4.2e-5}{J/m^3}$；每米光束的 $P/c = \qty{3.3e-11}{J}$，
在 \qty{30}{cm} 的距离内共有 $4.7 \times 10^5$ 个波长。

\textbf{4.} $\langle\Pi\rangle = I$ 沿 $z$ 方向，按 $\cos^2$ 规律以 $2f$ 振荡；
黑体上为 $P/c = \qty{3.3e-11}{N}$，镜子上为 \qty{6.7e-11}{N}。

\textbf{5.} $E_0/5 \times 10^{11} = 6 \times 10^{-9}$；$B_0$ 为地球
磁场的五分之一。

\textbf{6.} $I/c = \qty{4.2e-5}{Pa}$。

\textbf{7.} $v_0 = eE_0/m\omega = \qty{0.18}{m/s}$；$F_m/F_e = v_0B_0/E_0 = v_0/c = 6 \times
10^{-10}$：慢速电荷所受的磁力可以忽略。

\textbf{8.} $\vect E = E_0[\cos(\omega t - kz)\vect e_x \pm \sin(\omega t - kz)\vect e_y]$，其中
$E_0 = \qty{2.2}{kV/m}$（功率相同），$\vect B = \vect e_z\wedge\vect E/c$：
强度不变，但 $\Pi = \varepsilon_0cE_0^2$ 现在\emph{恒定}
不随时间变化——能量流不再脉动。

\textbf{9.} $10\cos^230^\circ = \qty{7.5}{mW}$；$0$；$7.5 \times \cos^445^\circ = \qty{1.9}{mW}$。

\textbf{10.} 亮：$\tfrac12$；暗：$0$。

\textbf{11.} 经过第一个起偏器后，两坐标轴上的分量相等；
经过延迟片后，两分量之间有相位差 $\delta$；在正交
方向上的投影 $\propto(1 - \eu^{\iu\delta})/2$：$I \propto \sin^2(\delta/2)$。

\textbf{12.} 暗；强度减半。

\textbf{13.} \omterm{def:b2:plane-waves-polarization:states}{圆偏振}的输出光经偏振太阳镜观察
时，无论头部转到何种角度，亮度都相同。

\textbf{14.} $\tfrac12 \times 0.85 \times 0.85 \times \tfrac13 = 12\%$；第一个起偏器
（吸收一半）和颜色滤光片（吸收三分之二）。

\textbf{15.} 亮 $0.36$，暗 $\tfrac12 \times 10^{-3}$：对比度 $\approx 700$。

\textbf{16.} \omterm{def:b2:plane-waves-polarization:states}{圆偏振}光：无论偏振片如何转动，
对偏振敏感的分束器接收到的功率
都相同。

\textbf{17.} \omterm{def:b2:plane-waves-polarization:states}{圆偏振}光的旋向在
头部倾斜时保持不变；\omterm{def:b2:plane-waves-polarization:states}{线偏振}眼镜倾斜 \qty{20}{\degree} 后，有 $\sin^220^\circ = 12\%$ 漏入
错误的那只眼睛。

\textbf{18.} $F = 2IS/c = \qty{9.1}{mN}$；$a = \qty{1.8e-4}{m/s^2}$；太阳
引力的 $3\%$。

\textbf{19.} 截获功率 $\propto\cos\theta$，反射引起的动量变化
沿法线方向为 $2p\cos\theta$：$F = (2IS/c)\cos^2\theta$，沿法线方向。

\textbf{20.} $F = 9.1 \times \cos^235^\circ = \qty{6.1}{mN}$，切向分量 $F\sin35^\circ
= \qty{3.5}{mN}$：$a_t = \qty{7e-5}{m/s^2}$，每年改变 \qty{2.2}{km/s}——
相当于轨道速率的 $7\%$。

\textbf{21.} $(1/0.39)^2 = 6.6$；引力同样按此比例变化。

\textbf{22.} $0.1 \times 1360 = 2\sigma T^4$：$T = \qty{190}{K}$——
镜子仍然是冷的。

\textbf{23.} $1000/1.8 \times 10^{-4} = \qty{5.6e6}{s}$，两个月；黑色帆
所受的力减半。

\textbf{24.} $1.36 \times 10^6/3.2 \times 10^{-19} = \qty{4.3e21}{}$ 每秒光子数，
每个光子的 $hf/c = \qty{1.1e-27}{kg.m/s}$；反射时 $2Np = \qty{9.1}{mN}$。

\textbf{25.} 能量（$u$、$I$）：光束挡板和光电二极管。
动量（$I/c$）：帆。偏振（\omterm{prop:b2:plane-waves-polarization:malus}{马吕斯定律}、偏振片）：
屏幕及其起偏器。
\end{solution}
